% New mean-one degree-dispersion crossover; all claims are within conditional DMFT.
\section{Bias, curvature and a disorder-induced critical crossover}
The three-halves exponent is selected by nonzero local curvature, rather than
by heterogeneity alone. Define the stationary response gap
$\delta=1-g^2A(q)$ and the critical moments
$A=\E\phi'(\mu)^2$, $B=\E\phi''(\mu)^2$ and
$C=\E\phi'''(\mu)^2$. Price expansion and the first integral give
\begin{align*}
 K''&=\delta K-\tfrac12g^2BK^2-\tfrac16g^2CK^3+\cdots,\\
 (K')^2&=\delta K^2-\tfrac13g^2BK^3-\tfrac1{12}g^2CK^4+\cdots,\\
 \delta&=\tfrac13g^2Bq+\tfrac1{12}g^2Cq^2+\cdots,\qquad q=K(0).
\end{align*}
For positive error-function rates and $\mu\sim\mathcal N(h,\Delta)$,
\[
 \frac BA=\frac{\Delta}{1+2\Delta}
              +\frac{h^2}{(1+2\Delta)^2}.
\]
Thus $B=0$ requires both $h=0$ and $\Delta=0$. A homogeneous biased operating
point already has $B>0$; a symmetric but spread distribution also has $B>0$.

For fixed degree $\chi=0$ and prescribed mean rate $m$, put
$a=\Phi^{-1}(m/2)$ and $r=g/g_c-1$. The homogeneous law has
$\Delta=0$, $h=a\sqrt{1+q}$ and $A(q)=\phi'(a)^2/(1+q)$. Therefore
$\delta=q-2r+\cdots$. Combining with the first integral gives, for $a^2<3$,
\[
 q\sim\frac6{3-a^2}r,\qquad
 \delta\sim\frac{2a^2}{3-a^2}r.
\]
At the benchmark $m=.3$ this is a three-halves case even without static
heterogeneity. At $m=1$ the quadratic term vanishes, $q\sim2r$ and
$\delta\sim r^2/3$, recovering the homogeneous odd cubic law.

\subsection{A uniform small-dispersion expansion at mean rate one}
For $m=1$, write $d=\Delta$, $s=1+d+q$ and
$\lambda=2g^2/\pi$. The exact static and energy equations are
\begin{align*}
 d&=\lambda\{\tfrac\pi2\chi+\arcsin(d/s)\},\\
 1&=2\lambda\int_0^1
       \frac{1-y}{\sqrt{s^2-(d+qy)^2}}\,\dd y.
\end{align*}
Expanding the energy integral gives
\[
 \lambda=1+d+q-\frac{d^2}2-\frac{dq}3-\frac{q^2}{12}
           +O((d+q)^3),\qquad
 \delta=\frac{dq}3+\frac{q^2}{12}+O(q(d+q)^2).
\]
Substitution into the static equation yields
\[
 \lambda\arcsin(d/s)
 =d-\frac{d^3}3-\frac{d^2q}3-\frac{dq^2}{12}+\cdots,
\]
and consequently
\[
 \boxed{d(2d+q)^2=6\pi\chi+O((d+q)^4+\chi(d+q)).}
\]
In particular the static \emph{variance} obeys
$\Delta_c\sim(3\pi\chi/2)^{1/3}$; the standard deviation has exponent $1/6$.
Let $d_c=(3\pi\chi/2)^{1/3}$, $Z=d/d_c$, $Q=q/d_c$ and
$r=g/g_c(\chi)-1$. The leading uniform curve is
\[
 Z(2Z+Q)^2=4,\qquad
 \frac r{d_c}=\frac{Z+Q-1}2,\qquad
 \frac\delta{d_c^2}=\frac{ZQ}3+\frac{Q^2}{12}.
\]
For $Q\ll1$, $Z=1-Q/3+\cdots$, hence $q\sim3r$ and $\delta\sim d_cr$.
For $Q\gg1$ with $q$ still small, $Z\sim4/Q^2$, $d\sim6\pi\chi/q^2$,
$q\sim2r$ and $\delta\sim r^2/3$. Thus the inner three-halves region has
width $r_\chi\asymp\chi^{1/3}$ and crosses to the odd homogeneous cubic law.
The two limits $\chi\to0$ and $r\to0$ do not commute in this exponent.

\subsection{Kernel and dimension through the crossover}
The mixed normal form has first integral
$(K')^2=K^2(\delta-\beta K-\gamma K^2)$, with
$\beta=d/3$, $\gamma=1/12$ and $\delta=\beta q+\gamma q^2$. Setting $U=1/K$
gives $U'^2=\delta U^2-\beta U-\gamma$. Its homoclinic solution is
\[
 K(\tau)=\frac{2\delta}
 {\beta+\sqrt{\beta^2+4\gamma\delta}\cosh(\sqrt\delta\tau)}.
\]
Writing $a=4d/(4d+q)$ and $s=\sqrt\delta\tau$,
\[
 k_a(s)=K/q=\frac2{a+(2-a)\cosh s}.
\]
This goes from $\operatorname{sech}^2(s/2)$ at $a=1$ to
$\operatorname{sech}(s)$ at $a=0$. For
$\vartheta=\arccos[a/(2-a)]$ its Fourier transform is
\[
 f_a(x)=\frac{4\pi\sinh(\vartheta x)}
 {(2-a)\sin\vartheta\sinh(\pi x)},
\]
with the continuous endpoint limit $f_1=4\pi x/\sinh(\pi x)$.
The narrow pair resonance then gives
\[
 \frac{\mathrm{PR}}N\sim\frac{A^2}{A_4}
       \frac{\delta^{3/2}}{\mathcal C(a)},\qquad
 \mathcal C(a)=\frac1{4\pi}\int_{-\infty}^\infty
                    \frac{f_a(x)^2}{1+x^2}\,\dd x.
\]
The endpoints are $\mathcal C(0)=\pi^2/12$ and
$\mathcal C(1)=4\pi^2/3-12$. Since $A^2/A_4\to1$ in the joint mean-one
small-dispersion limit,
\[
 \frac{\mathrm{PR}/N}{\chi}\longrightarrow
 \frac{3\pi}{2}\frac{(ZQ/3+Q^2/12)^{3/2}}
 {\mathcal C(4Z/(4Z+Q))}.
\]
This is a parameter-free scaling curve within the conditional theory.

\subsection{Visibility and scope}
The extensive result requires $N(\mathrm{PR}/N)\gg1$. In the inner region
this implies
\[
 \frac{\mathcal C(1)^{2/3}}{d_cN^{2/3}}\ll r\ll d_c.
\]
A nonempty window therefore requires $N\chi\gg1$, up to constants. This
condition is necessary, not sufficient; sample-dependent thresholds and
recording errors may impose stronger conditions. At crossover
$\tau_{\rm pair}\asymp\chi^{-2/3}$. Nearer onset it grows as $(d_cr)^{-1}$,
whereas the local timescale grows as $(d_cr)^{-1/2}$.
The diagonal-pair-lag time is not by itself a universal recording criterion.
The block-centering energy contraction is $\psi(0,\tau)$, whose integrated
time is instead
\[
 \tau_E\sim\frac{f_a(0)^2}{4\mathcal C(a)}\delta^{-1/2}.
\]
This equals $4\delta^{-1/2}/\mathcal C(1)$ at the quadratic endpoint and
$3\delta^{-1/2}$ at the cubic endpoint. Actual estimator variance and gap
conditions additionally require the corresponding connected four-time
contraction; the slower $\delta^{-1}$ scale can be invoked only after showing
that the chosen observable couples to that sector.

The exact conditional-DMFT branch and resonant pair checks are retained in
\texttt{research\_round/critical}. At $\chi=10^{-9}$ and $Q=.1,1,10$, pair
integrals divided by the mixed-shape prediction using the actual gap and
variance are $1.000159,1.001054,1.004108$. These are deterministic tests
inside the closure, not network exponent measurements. The background
three-halves claim must retain its transversality and local-curvature domain;
the exponent alone does not establish novelty or a special role for balance.
